\documentclass[11pt,a4paper]{article}

\usepackage[margin=1in]{geometry}
\usepackage[T1]{fontenc}
\usepackage[utf8]{inputenc}
\usepackage{lmodern}
\usepackage{amsmath,amssymb,amsthm}
\usepackage{enumitem}
\usepackage{hyperref}

\title{Contextual Bandit Formulation for Fintech Churn Retention}
\author{}
\date{}

\begin{document}

\maketitle

\section{Problem Setting}

We model churn-retention as a stochastic contextual bandit over \(T\) decision rounds.  
At each round \(t=1,\dots,T\):

\begin{enumerate}[label=\arabic*.]
    \item A customer context \(x_t \in \mathcal{X} \subseteq \mathbb{R}^d\) is observed.
    \item A feasible set of retention actions \(A_t \subseteq A=\{1,\dots,K\}\) is available.
    \item The decision policy selects an action \(a_t \in A_t\).
    \item A reward \(r_t \in \mathbb{R}\) is observed only for the chosen action.
\end{enumerate}

The objective is to learn a policy \(\pi(a \mid x)\) maximizing expected cumulative reward:
\[
\max_{\pi} \; \mathbb{E}_{\pi}\left[\sum_{t=1}^T r_t\right].
\]

\section{Fintech Mapping}

\subsection*{Context}

The context vector \(x_t\) summarizes the state of customer \(i_t\) at time \(t\), for example:
\begin{itemize}
    \item demographic and segment variables,
    \item behavioral variables such as recency, frequency, and monetary value,
    \item engagement signals such as app usage and support contacts,
    \item churn-risk and customer lifetime value estimates,
    \item previous offers and responses.
\end{itemize}

\subsection*{Actions}

Each action corresponds to a retention measure:
\[
A=\left\{
a^{(0)}\text{ no action},\;
a^{(1)}\text{ message},\;
a^{(2)}\text{ fee waiver},\;
a^{(3)}\text{ cashback},\;
a^{(4)}\text{ rate boost},\;
a^{(5)}\text{ human outreach}
\right\}.
\]

The feasible subset \(A_t\) may depend on eligibility, compliance, and budget constraints.

\subsection*{Reward}

Let \(V_H(x_t,a)\) denote the expected customer value over a fixed horizon \(H\) if action \(a\) is applied in context \(x_t\).  
Using \(a^{(0)}\) (no action) as the baseline, define the incremental value of action \(a\) by
\[
\Delta V_H(x_t,a) = V_H(x_t,a) - V_H(x_t,a^{(0)}).
\]

We then define the reward as the \emph{incremental net value}
\[
r_t = \Delta V_H(x_t,a_t) - c(a_t),
\]
where \(c(a_t)\geq 0\) is the cost of the selected retention action.

Equivalently,
\[
r_t = V_H(x_t,a_t) - V_H(x_t,a^{(0)}) - c(a_t).
\]

This definition is important: the policy does not maximize raw future revenue, but only the \emph{additional} value created by intervening relative to doing nothing. Hence, for a customer who is already valuable and unlikely to churn, the optimal decision may be \(a^{(0)}\).

\section{Statistical Model}

Let \(\phi(x_t,a)\in\mathbb{R}^m\) be a joint feature map for context-action pairs. A practical choice is
\[
\phi(x_t,a)=\bigl[x_t;\; \mathrm{onehot}(a);\; x_t \odot \mathrm{onehot}(a)\bigr],
\]
so that the model captures both customer effects and action-specific interactions.

A linear contextual bandit assumes
\[
\mathbb{E}[r_t \mid x_t, a_t=a] = \phi(x_t,a)^\top \theta^\star,
\]
for an unknown parameter vector \(\theta^\star \in \mathbb{R}^m\).

Under this specification, the model directly learns the expected incremental net value of each action for each customer context.

\section{Recommended Learning Rule}

A suitable algorithm is linear contextual Thompson Sampling. We assume
\[
r_t = \phi(x_t,a_t)^\top \theta^\star + \varepsilon_t,
\qquad
\varepsilon_t \sim \mathcal{N}(0,\sigma^2),
\]
with Gaussian prior
\[
\theta^\star \sim \mathcal{N}(\mu_1,\Sigma_1).
\]

At round \(t\), given the posterior
\[
\theta^\star \mid \mathcal{D}_{t-1} \sim \mathcal{N}(\mu_t,\Sigma_t),
\]
the algorithm proceeds as follows:

\begin{enumerate}[label=\arabic*.]
    \item Sample
    \[
    \tilde{\theta}_t \sim \mathcal{N}(\mu_t,\Sigma_t).
    \]
    \item Select
    \[
    a_t
    =
    \arg\max_{a \in A_t}
    \phi(x_t,a)^\top \tilde{\theta}_t.
    \]
    \item Observe the reward \(r_t\).
    \item Update the posterior parameters:
    \[
    \Sigma_{t+1}^{-1}
    =
    \Sigma_t^{-1}
    +
    \frac{1}{\sigma^2}\phi_t\phi_t^\top,
    \]
    \[
    \mu_{t+1}
    =
    \Sigma_{t+1}
    \left(
    \Sigma_t^{-1}\mu_t
    +
    \frac{1}{\sigma^2}\phi_t r_t
    \right),
    \]
    where \(\phi_t := \phi(x_t,a_t)\).
\end{enumerate}

Equivalently, one may define
\[
V_t = \Sigma_t^{-1},
\qquad
b_t = V_t \mu_t,
\]
so that the update becomes
\[
V_{t+1} = V_t + \frac{1}{\sigma^2}\phi_t\phi_t^\top,
\qquad
b_{t+1} = b_t + \frac{1}{\sigma^2}\phi_t r_t,
\]
and
\[
\mu_{t+1} = V_{t+1}^{-1} b_{t+1}.
\]

Thus, after each observed reward, the posterior mean and covariance are
updated using the chosen context-action feature vector \(\phi_t\), and the
next action is sampled according to the updated posterior.

\section{Recommended Learning Rule}

A suitable algorithm is contextual Thompson Sampling with linear rewards.
At each round:

\begin{enumerate}[label=\arabic*.]
    \item Maintain a posterior distribution over \(\theta^\star\).
    \item Sample \(\tilde{\theta}_t\) from the posterior.
    \item Select
    \[
    a_t
    =
    \arg\max_{a \in A_t}
    \phi(x_t,a)^\top \tilde{\theta}_t.
    \]
    \item Observe \(r_t\) and update the posterior.
\end{enumerate}

Since \(a^{(0)}\) is part of the action set, the algorithm can explicitly learn when the best intervention is to do nothing.

\end{document}